\begin{claim}{D.36}{Proved}{The global multiplier of a positive transposition}
Let $B$ be the full three-cusp loop-outer group, and let
$g\in B$ be an involution whose $2$-adic multiplier is
$\lambda_2(g)=1$.  Its puncture permutation is a
transposition, and its common multiplier is the global
element $1\in\widehat\Z^\times$.

Consequently a finite subgroup of $B$ whose map (D.32a)
has image $S_3\times C_2$ has the usual global orientation
character.  The same is true when that image is a
$V_4$ containing the pure involution, or is the positively
oriented $S_3\times\{1\}$.  These statements identify
the actual multiplier homomorphism; they do not assert
conjugacy of those finite actions.
\end{claim}
\begin{proof}
We first record the genus-zero extension of the actual
automorphism calculation already proved in D.32--D.33.
Let $F_n=\widehat\pi_{0,n}$ with $n\geq3$, and let
$\beta$ be a pure automorphism with $\beta^p=1$ and
$p$-adic peripheral multiplier one.  Then $\beta=1$.
Indeed $F_n$ is free of rank $r=n-1$, and its cusp
classes span its abelianization.  Its induced pro-$p$
automorphism has identity abelianization and finite
order, so the $\Z_p$ Magnus argument of D.33 makes
it the identity.  The ring calculation (D.33c)--(D.33g)
gives one conjugacy class of sections of
$F_n\rtimes_\beta C_p\to C_p$, and its actual fixed
subgroup $C$ has
\[
                       \dim H^1(C,\F_p)=n-1.               \tag{D.36a}
\]
The norm calculation of D.28 applies to every cusp:
its image is a primitive element of the free pro-$p$
group, the first coordinate of its normalizer lift
centralizes that element, and the lift's $p$th power
has $p$-coordinate divisible by $p$.  Coordinatewise
division of the norm therefore supplies a section in
each peripheral normalizer.  Uniqueness of the
section class places conjugates of all $n$ full
peripheral $\Z_p$ subgroups inside the same $C$.

The finite-cover rank proof of D.32 now applies with
all $n$ cusps.  If $C$ is open, its first-cohomology
dimension is $1+[F_n:C](n-2)$, which equals $n-1$
only at index one.  If it has infinite index and
finite $p$-part of that index, the same exact
restriction colimit and eventual prime-to-$p$
injections give unbounded first-cohomology dimension.
Otherwise choose an open $U\geq C$ of $p$-divisible
index.  The $n$ selected boundary lifts all have
degrees prime to $p$ and lie over different original
boundaries.  They cannot be all the boundaries:
a single lift over any original boundary would
have degree $[F_n:U]$.  With at least $n+1$ boundaries,
the $n$ selected boundary homology classes are
independent.  Their dual characters restrict to
$C$ with the nonzero multipliers $d_i^{-1}\bmod p$
computed in (D.32d).  Thus $\dim H^1(C,\F_p)\geq n$,
contradicting (D.36a).  Hence $C=F_n$ and $\beta=1$.
This proof concerns an actual finite-order automorphism;
it does not presume a lift for arbitrary finite outer
actions in larger rank.

Return to $g$.  D.32 excludes a nonidentity pure
involution with $\lambda_2=1$, so its permutation
must be a transposition.  Relabel the original cusp
generators as $x,y,z=(xy)^{-1}$ so that $x$ is fixed
and $y,z$ are exchanged.  D.27 provides an actual
order-two lift $\alpha$ of this outer action.
Its peripheral multiplier satisfies $\lambda^2=1$,
so all its prime coordinates are $1$ or $-1$.

Suppose $\lambda_q=-1$ for an odd prime $q$, and use
the exact cyclic quotient
\[
 \chi:F\longrightarrow C_q,\qquad
       \chi(x)=0,\quad\chi(y)=1,\quad\chi(z)=-1,
                   \qquad N=\ker\chi.                    \tag{D.36b}
\]
The puncture permutation and the multiplier give
$\chi\alpha=\chi$: on $y$ the exponent $-1$ at
$q$ is multiplied by $\chi(z)=-1$, and similarly
for $z$.  Thus $N$ is invariant and the induced
action on its cyclic deck quotient is the identity.
The degree-$q$ cover is a sphere with $q+2$ cusps.
There are $q$ degree-one cusps over $x$, and one
degree-$q$ cusp over each of $y,z$.

Here are the exact covering and filling generators.
Finite-coset Schreier rewriting with coset representatives
$1,y,\ldots,y^{q-1}$ gives the free basis of $N$
\[
 a_i=y^i x y^{-i}\quad(0\leq i<q),\qquad b=y^q.
\]
The remaining boundary generator is
\[
 d=z^q=b^{-1}a_{q-1}^{-1}\cdots a_0^{-1},
                     \qquad a_0\cdots a_{q-1}bd=1.        \tag{D.36c}
\]
For example the latter relation is the actual free-word
identity
$(xy)^q=a_0\cdots a_{q-1}y^q$.
Completing the discrete Schreier comparison identifies
this presented free profinite group with the specified
open subgroup $N$.

The lift $\alpha$ fixes the conjugacy class of every
$a_i$ inside $N$.  To see this without choosing an
unproved normalization, write
$\alpha(x)=h^{-1}x^\lambda h$ for some $h\in F$.
Since $\chi\alpha(y)=1$, its action on the indices
of these cusp classes is the translation
$i\mapsto i-\chi(h)$ in $C_q$.  Its square is the
identity, so $2\chi(h)=0$ and $\chi(h)=0$.
Explicitly
\[
 k_i=\alpha(y)^i h^{-1}y^{-i}\in N,\qquad
                \alpha(a_i)=k_i a_i^\lambda k_i^{-1}.     \tag{D.36d}
\]
The two unique boundary classes over $y,z$ are
exchanged, with multiplier $\lambda$.

Fill exactly these two cusps.  The map is the quotient
by the closed normal subgroup of $N$ generated by $b,d$:
\[
 N\longrightarrow
 K=N/\overline{\langle\!\langle b,d\rangle\!\rangle}_N
   =\widehat{\langle a_0,\ldots,a_{q-1}\mid
                             a_0\cdots a_{q-1}=1\rangle}
   \cong\widehat\pi_{0,q}.                               \tag{D.36e}
\]
Both presentation maps are explicit: kill $b$ and
send the remaining $a_i$ to the named peripheral
generators; the inverse identifies $a_{q-1}$ with
$(a_0\cdots a_{q-2})^{-1}$.  These substitutions
verify both universal properties and compositions.
Since $\alpha$ exchanges the normal generators'
peripheral classes, the closed normal subgroup
in (D.36e) is invariant.  Consequently $\alpha$
induces an actual automorphism $\bar\alpha$ of
$K$ with $\bar\alpha^2=1$.  By (D.36d) it is
pure, with the same multiplier $\lambda$.

The proved genus-zero calculation at $p=2$ applies:
$q\geq3$ and $\lambda_2=1$, so $\bar\alpha=1$.
But its $q$-adic multiplier on each surviving
nontrivial cusp is $-1$, which is not the
identity on the cusp's $\Z_q$ subgroup.  Equivalently,
it acts as $-I$ on the nonzero $q$-adic abelianization
of the genus-zero group $K$.  This contradiction
excludes every such $q$.  Together with
$\lambda_2=1$ it proves $\lambda=1$ globally.

For the finite-group consequences, use the faithful
map (D.32a).  In a $V_4$ with pure kernel, the
pure involution has multiplier $-1$ by D.34,
and the positive transposition has multiplier
$1$ by the present proof.  These generate it.
The positively oriented $S_3$ is generated by
its positive transpositions.  Finally
$S_3\times C_2$ is generated by those
transpositions and its pure involution.
Multiplicativity of the common multiplier proves
each asserted character identity.

The program \texttt{code/part\_c/positive\_transposition\_d36.py}
checks (D.36c) as an exact free-word identity,
the finite cyclic cusp permutations and their
involutive translations, and the two presentation
substitutions in (D.36e) for several odd primes.
The proof above supplies the all-prime assertion.
\end{proof}
